\section{Lecture 10, Part 0: Model Evaluation}

\begin{frame}
\vfill
\centering
{\color{RLBlue}\large Lecture 10}

\vspace{0.3em}
{\color{RLBlue}\Large\bfseries Part 0}

\vspace{0.6em}
{\Large Model Evaluation}
\vfill
\end{frame}

\begin{frame}{Learning Goals}
\begin{itemize}
  \item Say what a training, validation and test set are each for, and why a test set is spent the moment you look at it.
  \pause
  \item Read the bias--variance tradeoff off a measured error curve, and explain why more capacity eventually hurts.
  \pause
  \item Put an honest interval around a reported metric using cross-validation and the bootstrap.
\end{itemize}
\end{frame}

% ─── Train / validation / test ──────────────────────────────────────────────
\begin{frame}{Why Do We Need a Test Set At All?}
\begin{itemize}
  \item \textbf{Why:} we never care how well a rule describes the data we already have. We care how it does on a day we have not seen. Those two numbers can be wildly different.
  \pause
  \item \textbf{Everyday example:} a student works through ten past papers until they can reproduce every answer from memory. Their score on those ten papers is $100\%$. It says nothing about tomorrow's exam, because the exam is the only paper they have not memorised.
  \pause
  \item \textbf{Numerical example:} the coffee shop fits a curve through $20$ days of iced-coffee sales. A wiggly enough curve scores an error of $1.5$ cups on those $20$ days -- and $27.7$ cups on $40$ fresh days.
  \pause
  \item The error on data you fitted is not an estimate of anything. It is a measure of how hard you tried.
\end{itemize}
\end{frame}

\begin{frame}{Training, Validation and Test}
\begin{rldefinition}
\textbf{Training set} -- the data the rule is fitted on. \emph{The past papers you practise with.}

\smallskip
\textbf{Validation set} -- held-out data used to \emph{choose} between candidate rules: how wiggly a curve, how deep a tree. \emph{The mock exam that tells you which revision method is working.}

\smallskip
\textbf{Test set} -- held out and untouched until the very end, used once to report how well the chosen rule does. \emph{The real exam.}
\end{rldefinition}
\pause
\begin{itemize}
  \item The split matters because each set answers a different question. Choosing on the test set turns it into a validation set, and then nothing is left to report with.
\end{itemize}
\end{frame}

\begin{frame}{Real Event: Fitting Twenty Days of Sales}
\begin{columns}[T]
\column{0.50\textwidth}
\small
\begin{itemize}
  \item The shop records the day's high temperature and iced coffees sold, for $20$ days. Sales stay flat while it is cool, then climb.
  \pause
  \item \textbf{Degree $1$} -- a straight line. It cannot bend, so it misses the flat stretch and the top: error $2.84$ training, $3.33$ test.
  \pause
  \item \textbf{Degree $3$} -- follows the real shape: $2.03$ and $2.65$.
  \pause
  \item \textbf{Degree $9$} -- passes near every dot, even predicting \emph{negative} sales near $21^\circ$C to get there: $1.48$ training, $27.7$ test.
\end{itemize}
\column{0.48\textwidth}
\centering
\includegraphics[width=\linewidth]{figures/evaluation/evaluation_polynomial_fits.pdf}
\end{columns}
\end{frame}

% ─── Bias-variance ──────────────────────────────────────────────────────────
\begin{frame}{Why Does More Capacity Stop Helping?}
\begin{itemize}
  \item \textbf{Why:} a fit can be wrong for two opposite reasons, and fixing one makes the other worse.
  \pause
  \item \textbf{Too rigid} -- a straight line cannot bend to the truth, however much data it gets. It is wrong the same way every time: that is \textbf{bias}.
  \pause
  \item \textbf{Too flexible} -- a degree-$9$ curve chases the noise in whichever $20$ days it saw. Show it $20$ different days and it comes out a different shape: that is \textbf{variance}.
  \pause
  \item \textbf{And some error is nobody's fault:} two days at the same temperature still sell different amounts.
\end{itemize}
\end{frame}

\begin{frame}{The Bias--Variance Tradeoff}
\begin{rldefinition}
The error a rule makes on fresh data splits into three parts:
\[
\underbrace{\text{bias}^2}_{\text{too rigid to fit the truth}}
\; + \;
\underbrace{\text{variance}}_{\text{changes too much with the sample}}
\; + \;
\underbrace{\text{noise}}_{\text{irreducible}}
\]
\emph{Everyday:} the student who memorised ten papers has low bias and huge variance -- ask a question they did not memorise and the answer is anyone's guess. The student who answers ``it depends'' to everything has low variance and huge bias.
\end{rldefinition}
\pause
\begin{itemize}
  \item Adding capacity trades bias down and variance up. There is a point where the trade stops paying, and only held-out data can tell you where it is.
\end{itemize}
\end{frame}

\begin{frame}{Where Does the Test Error Turn Around?}
\begin{columns}[T]
\column{0.50\textwidth}
\small
\begin{itemize}
  \item Same $20$ training days, every degree from $1$ to $10$, error measured on the $40$ test days.
  \pause
  \item \textbf{Training error only ever falls} -- $2.84$ down to $1.47$. It has to: more capacity can always fit the same points better.
  \pause
  \item \textbf{Test error falls, bottoms out at degree $3$ ($2.65$), then climbs} -- $4.8$ at degree $7$, then $22$, $28$, $34$. It wobbles on the way: $40$ test days is not many.
  \pause
  \item \textbf{Read it:} the gap between the curves \emph{is} the overfitting. Watching only the green line, degree $10$ looks like your best model ever.
\end{itemize}
\column{0.48\textwidth}
\centering
\includegraphics[width=\linewidth]{figures/evaluation/evaluation_bias_variance.pdf}
\end{columns}
\end{frame}

\begin{frame}{Why This Matters: Capacity Is Not Free}
\begin{itemize}
  \item \textbf{The promise, cashed:} every model you meet from here on has a capacity dial -- the depth of a decision tree, the number of neighbours in k-NN, the width of a network, the strength of a penalty term.
  \pause
  \item \textbf{Repeat this one:} \emph{training error always falls with capacity, so it can never tell you when to stop.} Only held-out data can.
  \pause
  \item \textbf{Module 1, Lecture 5} prunes decision trees for exactly this reason: an unpruned tree can put every training point in its own leaf, which is degree $10$ wearing a different hat.
  \pause
  \item \textbf{Module 1, Lecture 14} returns to this curve by name, and the whole of \textbf{regularization} (Module 1, Lecture 12) is machinery for holding variance down without giving up the flexibility.
\end{itemize}
\end{frame}

% ─── Cross-validation ───────────────────────────────────────────────────────
\begin{frame}{Why Do We Need Cross-Validation?}
\begin{itemize}
  \item \textbf{Why:} holding out a validation set costs you data twice -- the held-out days cannot be trained on, and the score from a handful of days is itself noisy.
  \pause
  \item \textbf{$k$-fold cross-validation:} split the training data into $k$ equal parts. Train on $k-1$ of them, score on the one left out, and repeat until every part has been the held-out one exactly once. Average the $k$ scores.
  \pause
  \item Every observation is used for training $k-1$ times and for scoring once. Nothing is wasted, and you get $k$ scores instead of one -- so you can see how much the answer wobbles.
  \pause
  \item \textbf{The cost is arithmetic:} you fit the model $k$ times instead of once. With $k = 5$ that is usually a bargain; on a model that takes a week to train, it is not.
\end{itemize}
\end{frame}

\begin{frame}{Worked Example: Five Folds by Hand}
\begin{itemize}
  \item Take the shop's $20$ training days, shuffle them, and cut into $5$ folds of $4$ days. Fit the degree-$3$ curve five times, each time scoring on the fold left out:
\end{itemize}
\centering
\small
\begin{tabular}{l|ccccc}
Fold held out & $1$ & $2$ & $3$ & $4$ & $5$ \\
\hline
Error (cups)  & $2.60$ & $1.57$ & $2.19$ & $2.58$ & $3.66$ \\
\end{tabular}

\vspace{0.4em}
\raggedright
\begin{itemize}
  \item[] \pause
  \item Mean $= 2.52$ cups, standard deviation $= 0.76$. \pause
  \item \textbf{The spread is the point.} Had we held out only fold $2$ we would have reported $1.57$; had we held out only fold $5$, $3.66$. Either number alone would have been an accident of which days landed where.
\end{itemize}
\end{frame}

% ─── Uncertainty on a metric ────────────────────────────────────────────────
\begin{frame}{Why Do We Need an Interval on a Metric?}
\begin{itemize}
  \item \textbf{Why:} a test score is computed from a \emph{sample} of fresh data, so Lecture 5 applies: it is a random variable with a standard error.
  \pause
  \item \textbf{The problem:} $\sigma/\sqrt{n}$ applies to a mean. A median, an RMSE or an AUC is not a mean, and has no tidy formula.
  \pause
  \item \textbf{The bootstrap:} resample the test set \emph{with replacement} to the same size, recompute the metric, repeat thousands of times. The spread of those values estimates the spread of the metric.
  \pause
  \item It works because the test set is our best stand-in for the population -- drawing from it imitates drawing a fresh test set, the thing we cannot afford.
\end{itemize}
\end{frame}

\begin{frame}{How Precise Is That Error Number?}
\begin{columns}[T]
\column{0.50\textwidth}
\small
\begin{itemize}
  \item The degree-$3$ curve scores $\mathrm{RMSE} = 2.65$ cups on the $40$ test days. Resample those $40$ days $10{,}000$ times.
  \pause
  \item Middle $95\%$ of the resampled scores: $[2.03,\; 3.24]$ -- an interval more than a cup wide.
  \pause
  \item \textbf{Read it:} ``the error is $2.65$'' is really ``the error is somewhere around $2$ to $3.2$, and $40$ days is not enough to say better''.
  \pause
  \item A rival model scoring $2.60$ has not beaten this one. Their intervals overlap almost entirely.
\end{itemize}
\column{0.48\textwidth}
\centering
\includegraphics[width=\linewidth]{figures/evaluation/evaluation_bootstrap.pdf}
\end{columns}
\end{frame}

\begin{frame}{Why This Matters: $0.3\%$ Is Not a Difference}
\begin{itemize}
  \item \textbf{The promise, cashed:} leaderboards, papers and release notes all report a single accuracy. That number is an estimate, and it comes with a standard error whether or not anyone prints it.
  \pause
  \item \textbf{Numerical example:} a classifier at $90\%$ accuracy on a $1000$-example test set has $\mathrm{SE} = \sqrt{0.9 \times 0.1 / 1000} = 0.0095$, so a $95\%$ interval of roughly $\pm 1.9$ percentage points.
  \pause
  \item So $90.0\%$ against $90.3\%$ is not a result. The intervals sit almost exactly on top of each other, and rerunning with a different test split could easily swap the winner.
  \pause
  \item \textbf{Repeat this one:} \emph{halving the uncertainty on a metric needs four times the test data} -- the $\sqrt{n}$ of Lecture 5. At $10{,}000$ examples the interval is $\pm 0.6$ points, and only then is $0.3\%$ worth arguing about.
\end{itemize}
\end{frame}
